% SBI4GALEV 2026 summary white paper — generated by whitepaper.py
% Written locally by google/gemma-4-31b-it (alan); context never sent to external tools.
% Generated 2026-06-26T16:30:39
\documentclass[11pt]{article}
\usepackage[margin=1in]{geometry}
\usepackage{parskip}
\usepackage{enumitem}
\usepackage{hyperref}
\usepackage{titlesec}

\title{Simulation-Based Inference for Galaxy Evolution (SBI4GALEV) 2026: Summary White Paper}
\author{alan (automated rapporteur)}
\date{June 2026}

\begin{document}

\maketitle

\begin{abstract}
The SBI4GALEV 2026 conference convened researchers to discuss the integration of Simulation-Based Inference (SBI) and generative AI into the study of galaxy evolution. Over four days, 52 talks detailed a transition from introductory pedagogical frameworks to high-complexity implementations. The meeting highlighted a strategic shift toward amortized inference, differentiable forward modeling, and field-level analysis to meet the data demands of Stage-IV surveys. This paper synthesizes the core methodological advancements, the physical applications across cosmic scales, and the emerging trends in robustness and standardization.
\end{abstract}

\section{Overview of Topics}
The conference coverage was dominated by the intersection of statistical machine learning and astrophysical simulation. The primary focus remained on \textbf{Simulation-Based Inference (SBI) Core Methods}, specifically the deployment of Neural Posterior Estimation (NPE) and Neural Ratio Estimation (NRE) to bypass intractable likelihoods. This was closely followed by applications in \textbf{Galaxy Evolution \& Formation}, where researchers addressed quenching, mass assembly, and the galaxy-halo connection. \textbf{Generative AI \& Deep Learning} provided the architectural backbone for these efforts, with significant discussion on the utility of diffusion models and normalizing flows.

Medium-weight themes included \textbf{Cosmology \& Large Scale Structure} and \textbf{Forward Modeling}, often discussed as a unified pipeline for creating synthetic observations of the cosmic web. \textbf{Spectral Energy Distributions (SED)} and \textbf{Bayesian Methods} focused on the statistical rigor of spectral fitting and model evidence. Niche but critical topics included the physics of the \textbf{Interstellar and Circumgalactic Medium (ISM/CGM)}, \textbf{Active Galactic Nuclei (AGN)}, and \textbf{Differentiable Programming} using JAX. Finally, specialized applications were presented in \textbf{Multi-Messenger Astronomy} (LISA), \textbf{Galactic Archaeology} (Milky Way/LMC dynamics), and the \textbf{Epoch of Reionization}.

\section{Hot Subjects and Technical Dominance}

\subsection{Core SBI and Generative Architectures}
The most discussed subject was the evolution of SBI methodologies. A clear trend emerged moving beyond basic Normalizing Flows toward \textbf{Diffusion Models and Score-Based Generative AI}. Giuseppe Viterbo illustrated how score-matching allows for "crispier" posterior samples and flexible post-training constraints. Similarly, Grégoire Aufort demonstrated diffusion-based SED fitting by framing inference as an inpainting problem to handle missing photometric bands. Amortized inference remained central, with Tom and Niall Jeffries discussing the efficiency of NPE and the reframing of Bayesian evidence as a classification problem to avoid high-dimensional integrals.

\subsection{Galaxy Evolution and Population Synthesis}
Physical applications focused heavily on the galaxy-halo connection and mass assembly. Sinan Deger and Hiranya Peiris presented the \textit{pop-cosmos} framework, using diffusion models to learn the joint distribution of galaxy properties and calibrate redshift distributions for weak lensing. Christian Kragh Jespersen utilized Graph Neural Networks (GNNs) to demonstrate an equivalence between a galaxy's spatial environment and its assembly history. Other key contributions included Jonah Rose's work on the DREAMS project to disentangle sub-grid feedback from cosmic variance, and the use of SBI to study the impact of AGN on host systems (Thomas Bebbington).

\subsection{The Simulation Pipeline: Forward Modeling and Differentiability}
The integration of \textbf{Differentiable Programming} was a recurring highlight. Suchetha Cooray introduced \textit{Tengri}, a JAX-based framework for high-dimensional joint posterior inference, while Andrew Green presented \textit{L-Galaxies AD}, a fully auto-differentiable semi-analytic model. This shift enables gradient-based sampling (HMC/NUTS) to replace random-walk MCMC. Forward modeling efforts, such as the \textit{GalSBI} framework presented by Luca Tortorelli and Silvan Fischbacher, emphasized pixel-level analysis to transition from statistics-bounded to systematics-bounded regimes.

\subsection{Specialized Applications and Niche Domains}
The conference maintained breadth through several targeted applications:
\begin{itemize}
    \item \textbf{Gas and Environment:} Tanmay Singh applied SBI to infer multiphase CGM cloud properties, while Caterina Bracci used it to map nebular structures in the ISM.
    \item \textbf{Cosmology and LSS:} Alex Saoulis demonstrated multifidelity SBI to reduce the number of high-fidelity N-body simulations required for weak lensing.
    \item \textbf{Galactic Archaeology:} Richard Brooks used SBI to constrain the reflex motion of the Milky Way induced by the LMC, and Shichen Su presented fast stellar age inference.
    \item \textbf{High Energy/Multi-Messenger:} Philippa Cole explored SBI for EMRI signals in LISA, and Joey Braspenning recovered gas temperature distributions from X-ray spectra.
    \item \textbf{Reionization:} Timo Kist used normalizing flows to constrain the IGM neutral fraction during the Epoch of Reionization.
\end{itemize}

\section{Comparison with Previous Year}
Comparing the 2026 proceedings with the 2025 abstract booklet reveals a significant maturation of the field:
\begin{enumerate}
    \item \textbf{From Introduction to Implementation:} The 2025 program was heavily pedagogical, featuring "Introduction to SBI" and "Introduction to JAX." In 2026, these were replaced by specific, high-complexity applications, such as the LMC reflex motion and CGM cloud properties.
    \item \textbf{Architectural Shift:} While 2025 focused on cINNs and basic Normalizing Flows (e.g., Nils Candebat), 2026 saw a dominant shift toward Diffusion models, Flow-matching UNets (Ming-Shau Liu), and Foundation Models (Eduardo Albuquerque Hartmann).
    \item \textbf{Standardization:} 2026 introduced a formal "SBI Data Challenge" (proposed by Luca), signaling a move toward community benchmarking and validation that was absent in the previous year.
    \item \textbf{Differentiability:} The use of JAX evolved from a "brief introduction" (2025) to the core of new simulation codes like \textit{L-Galaxies AD} and \textit{Ceridwen} (2026).
\end{enumerate}

\section{Outlook for the Future}
The conference identified three primary frontiers for the coming years:
\begin{itemize}
    \item \textbf{Field-Level Inference:} A transition from summary statistics to pixel-level/field-level analysis to maximize information gain from Stage-IV surveys.
    \item \textbf{Robustness and Model Misspecification:} Increased focus on "Generalized Bayes" and MMD-based summaries to handle real-world data contamination, as exemplified by the \textit{OASIS} framework (Arya Farahi) and minimum-distance summaries (Sherman Khoo).
    \item \textbf{Interpretable AI:} The use of symbolic regression (Lucas Makinen, Deaglan Bartlett) and SHAP analysis (Shun-ya S. Uchida) to ensure that the "black box" of SBI reveals human-readable physical laws.
\end{itemize}

\end{document}
